\documentclass{jsarticle}
\usepackage{amsmath}
\usepackage{amssymb}

\begin{document}

\section{find\_fg\_4プログラムについて}

本プログラムは、指定された法 $P$ のもとでの有限環 $\mathbb{Z}/P\mathbb{Z}$ 上におけるアフィン変換群を基盤とし、特定の代数的制約およびグラフ上のサイクル制約を満たすパラメータ集合の完全探索を行うものである。探索の対象となる変数は $L$ 個（本実装では $L=12$）存在し、これらは前半部分（群 $F$、インデックス $0 \le i < L_H$）と後半部分（群 $G$、インデックス $L_H \le j < L$）に二分される。

\subsection{アフィン変換の定義と合成・逆元}
各インデックスに対応する変換を、写像 $f: \mathbb{Z}/P\mathbb{Z} \to \mathbb{Z}/P\mathbb{Z}$ とする。この写像は、変数 $a$ および $b$ を用いて以下のアフィン変換として定義される。
\[
f(x) \equiv a x + b \pmod P
\]
本系において逆写像が存在するための条件として、一次の係数 $a$ は $\gcd(a, P) = 1$ を満たすものに限定される。二つのアフィン変換 $f_1(x) \equiv a_1 x + b_1$ と $f_2(x) \equiv a_2 x + b_2$ の合成写像 $f_1 \circ f_2$ は、次式で与えられる。
\[
f_1(f_2(x)) \equiv a_1(a_2 x + b_2) + b_1 \equiv (a_1 a_2)x + (a_1 b_2 + b_1) \pmod P
\]
また、アフィン変換 $f(x)$ の逆写像 $f^{-1}(x)$ は、モジュラ逆数 $a^{-1}$ を用いて以下のように導出される。
\[
f^{-1}(x) \equiv a^{-1} x - a^{-1}b \pmod P
\]
プログラム内の関数 \texttt{composite} および \texttt{func\_inv} は、これらの代数構造を直接的に実装したものである。

\subsection{変換の可換性と線形合同式による空間縮小}
群 $F$ に属する変換 $F_i$ と、群 $G$ に属する変換 $G_j$ の間には、可換性（交換法則）に関する強い制約が存在する。$F_i \circ G_j = G_j \circ F_i$ が成り立つための条件は、それぞれの変換を展開することで以下のように示される。
\[
a_{F_i}(a_{G_j} x + b_{G_j}) + b_{F_i} \equiv a_{G_j}(a_{F_i} x + b_{F_i}) + b_{G_j} \pmod P
\]
両辺の一次の項は一致するため、定数項の比較により以下の線形合同式が得られる。
\[
(a_{G_j} - 1)b_{F_i} \equiv (a_{F_i} - 1)b_{G_j} \pmod P
\]
本アルゴリズムでは、探索の各段階において一方の変換パラメータが既知である性質を利用し、この線形合同式 $Ax \equiv B \pmod P$ を解くことで、他方の定数項 $b$ の取り得る候補を大幅に絞り込んでいる。ただし、特定のペア（インデックスの相対関係が特定の条件を満たす場合）においては「可換であってはならない（反交換的あるいは非可換である）」という例外的な制約が課される。その場合、上記合同式の解集合に属する $b$ を候補から明示的に除外する処理が行われている。

\subsection{サイクル非閉包条件と不動点の回避}
探索空間における最も厳密なグローバル制約は、グラフ上のサイクルに沿って定義された変換列に関する「非閉包条件」である。与えられたサイクル（巡回列）に沿って、各ノードにおけるアフィン変換およびその逆変換を順次合成して得られる最終的な変換を $f_{\text{cycle}}(x) \equiv A x + B \pmod P$ とする。

この巡回変換が「閉じる（元の状態にマッピングされる）」状態とは、変換が少なくとも一つの不動点を持つこと、すなわちある $x$ について $x \equiv A x + B \pmod P$ が成立することを意味する。式を整理すると以下のようになる。
\[
(A - 1)x \equiv B \pmod P
\]
この方程式が解 $x$ を持つための必要十分条件は、定数項 $B$ が法 $P$ と係数 $A - 1$ の最大公約数によって割り切れることである。
\[
B \equiv 0 \pmod{\gcd(A - 1, P)}
\]
プログラム内の \texttt{is\_closed} 関数はこの数理的性質を利用した不動点判定論理である。本系では「いずれのサイクルも不動点を持ってはならない（完全な閉包を構成しない）」ことが要求されているため、巡回変換の結果が上記の条件を満たした場合、そのパラメータの組み合わせを無効として棄却し、バックトラックを行う設計となっている。
\section{アフィン変換群に基づく探索の数理モデル（詳細展開）}

本プログラムの探索アルゴリズムは、有限環上のアフィン変換に関する代数的性質を巧みに利用し、探索空間を劇的に削減している。以下に、プログラム内で用いられている数式の詳細な導出と定式化を示す。

\subsection{有限環上のアフィン変換と可逆条件}
有限環 $\mathbb{Z}/P\mathbb{Z}$ 上の任意の要素 $x$ に対するアフィン変換 $f(x)$ は、傾き $a$ と切片 $b$ を用いて次のように定義される。
\[
f(x) \equiv a x + b \pmod P
\]
この写像が全単射（バイジェクション）となり、逆写像を持つための必要十分条件は、係数 $a$ と法 $P$ が互いに素であることである。
\[
\gcd(a, P) = 1
\]
プログラム内では、初期化段階においてこの条件を満たす $a$ のみが抽出され、探索候補として扱われている。

\subsection{変換の合成と逆写像の定式化}
探索中、グラフ上のパスに従って複数のアフィン変換が合成される。二つのアフィン変換 $f_1(x) \equiv a_1 x + b_1$ と $f_2(x) \equiv a_2 x + b_2$ の合成写像 $f_1 \circ f_2$ は次のように導出される。
\begin{align*}
(f_1 \circ f_2)(x) &= f_1(f_2(x)) \\
&\equiv a_1(a_2 x + b_2) + b_1 \pmod P \\
&\equiv (a_1 a_2)x + (a_1 b_2 + b_1) \pmod P
\end{align*}
これにより、合成写像もまた新たなパラメータ $a' = a_1 a_2$ および $b' = a_1 b_2 + b_1$ を持つアフィン変換として閉じていることがわかる。

また、ある変換 $f(x) \equiv a x + b$ に対する逆写像 $f^{-1}(y) = x$ は、方程式 $y \equiv a x + b \pmod P$ を $x$ について解くことで得られる。
\begin{align*}
a x &\equiv y - b \pmod P \\
x &\equiv a^{-1} (y - b) \pmod P \\
x &\equiv a^{-1} y - a^{-1} b \pmod P
\end{align*}
ここで $a^{-1}$ は法 $P$ における $a$ のモジュラ逆数（$\gcd(a, P)=1$ より一意に存在）である。これにより、逆写像のパラメータは $a_{\text{inv}} = a^{-1}$、$b_{\text{inv}} = -a^{-1} b$ として定式化される。

\subsection{可換性制約に伴う線形合同式の一般解法}
群 $F$ の変換 $f_F(x) \equiv a_F x + b_F$ と、群 $G$ の変換 $f_G(x) \equiv a_G x + b_G$ が可換（$f_F \circ f_G = f_G \circ f_F$）であるための条件は、以下の通りである。
\begin{align*}
a_F(a_G x + b_G) + b_F &\equiv a_G(a_F x + b_F) + b_G \pmod P \\
a_F a_G x + a_F b_G + b_F &\equiv a_G a_F x + a_G b_F + b_G \pmod P \\
a_F b_G + b_F &\equiv a_G b_F + b_G \pmod P \\
(a_G - 1)b_F &\equiv (a_F - 1)b_G \pmod P
\end{align*}
探索アルゴリズムにおいて、一方の変換パラメータ（例えば $G$ 側）と他方の傾き（$a_F$）が既知である場合、未知数 $x = b_F$ に関する一次合同式 $A x \equiv B \pmod P$ が構成される。
\[
A = a_G - 1, \quad B = (a_F - 1)b_G
\]
この線形合同式を解くため、まず最大公約数 $g = \gcd(A, P)$ を求める。解が存在する条件は $B \equiv 0 \pmod g$（すなわち $B$ が $g$ で割り切れること）である。この条件を満たすとき、元の式は以下のように $g$ で除算され、縮小された法 $M' = P/g$ のもとでの方程式に変換される。
\[
\frac{A}{g} x \equiv \frac{B}{g} \pmod{M'}
\]
ここで、$\gcd(A/g, M') = 1$ であるため、$(A/g)$ のモジュラ逆数が存在する。基本解 $x_0$ は次式で求まる。
\[
x_0 \equiv \left(\frac{A}{g}\right)^{-1} \frac{B}{g} \pmod{M'}
\]
法 $P$ における $x$ の解は全部で $g$ 個存在し、次のように表される。
\[
x_k \equiv x_0 + k M' \pmod P \quad (k = 0, 1, \dots, g-1)
\]
プログラム内の \texttt{solve\_linear\_congruence} 関数は、この厳密な数学的裏付けに基づいて解集合 $x_k$ を算出し、探索空間を効果的に刈り込んでいる。

\subsection{サイクル上の合成写像と不動点定理（非閉包条件）}
グラフ上の閉路（サイクル）に沿ってアフィン変換とその逆変換を反復的に合成した結果得られる最終的な巡回変換を $f_{\text{cycle}}(x) \equiv A x + B \pmod P$ とする。

この巡回変換が「系として閉じてしまう（元の状態に戻る）」というのは、数学的には変換が「不動点（Fixed Point）」を持つことと同値である。すなわち、ある入力 $x$ に対して $f_{\text{cycle}}(x) \equiv x \pmod P$ となる $x$ が存在するかどうかを検証する。
\begin{align*}
A x + B &\equiv x \pmod P \\
(A - 1)x &\equiv -B \pmod P
\end{align*}
定数の符号は法 $P$ においては吸収可能であるため、これは $(A - 1)x \equiv B \pmod P$ と本質的に同等である。

前項の線形合同式の解の存在条件と同様に、この方程式が解を持つ（不動点が存在する）ための必要十分条件は、定数項 $B$ が係数 $(A - 1)$ と法 $P$ の最大公約数 $g' = \gcd(A - 1, P)$ の倍数になっていることである。
\[
B \equiv 0 \pmod{\gcd(A - 1, P)}
\]
$A = 1$ の場合（恒等変換に近い状態）、$\gcd(0, P) = P$ となるため、条件は $B \equiv 0 \pmod P$（すなわち $B = 0$）に帰着する。プログラム内の \texttt{is\_closed} 関数はこの不動点定理の証明をそのまま論理式として実装しており、この合同式が成立した場合（不動点が存在しサイクルが閉じてしまった場合）にのみ、その探索パスを棄却している。
\end{document}